\section{Conclusion}
\label{sec:conclusion}
\enlargethispage{\baselineskip}

For the systems tested using \ac{fmri} responses to English-language stimuli, neither intervention route established a participant-general training benefit attributable to brain-response content.\evd{E008}\evd{E025}\evd{E026}\evd{E030}
The recorded-response route supported controlled brain predictivity, but did not demonstrate a participant-general pairing advantage across the tested objectives and variants.\evd{E002}\evd{E006}\evd{E008}\evd{E011}\evd{E013}\evd{E017}
WikiText assays in the synthetic-response route supported target uptake and retained-student movement at acceptable measured \ac{lm} quality.\evd{E016}\evd{E025}
However, the exact-target measurability criterion was not met, the text-derived auxiliary-control arm was non-identifying, and participant-level biological transfer was not demonstrated.\evd{E025}\evd{E026}\evd{E030}
The completed recorded-response utility evaluation did not support a brain-specific practical-utility claim.\evd{E009}

A controlled brain predictivity score alone does not establish a usable training signal.
Within the evidence standard used in this thesis, such a claim requires an ordered set of separately supported steps.
Before optimization, quality-aware headroom must be established and, for a synthetic target, exact-target measurability must also be established.
After optimization, target uptake and retained-student movement at acceptable measured \ac{lm} quality, brain-response-specific attribution, participant-level transfer to independently recorded responses, and the claimed utility endpoint must be evaluated separately.
For recorded-response training, a permuted-response control tests whether correct stimulus--response pairing matters.
For synthetic-response training, an adequately matched text-derived auxiliary-control arm tests whether any observed benefit is specific to brain-response content.
Seeds are technical repetitions, folds are data partitions, and voxels and \acp{roi} index brain-response measurements.
None can substitute for participants when the claim concerns people.
Keeping these steps separate shows what an intervention achieved, where the evidence chain stopped supporting the intended claim, and what evidence remains necessary before a controlled brain predictivity score can be treated as a useful training signal.
